\begin{claim}{F.17}{Proved}{The actual branch drilling and its nonfactorization}
Retain $Y,s_0,\zeta$ and the labelled map $\phi$ of F.12.
Write
\[
 v_u(s)=\left(\frac{u-s}{u+s}\right)^2\quad(u\in\mu_5),
 \qquad a(s)=\frac{\zeta-s}{\zeta+s},
 \qquad v_\zeta=a^2.
                                                               \tag{F.17.1}
\]
Thus $v_{\zeta^j}=v_{j+1}$ in the earlier notation.
For $u\in\mu_5\setminus\{\zeta\}$ put
\[
 A_u=u(2\zeta-u),\qquad B_u=\frac{\zeta u^2}{2u-\zeta},
 \qquad
 \Delta=\{s:s^2=A_u\text{ or }s^2=B_u
                    \text{ for one such }u\}.
                                                               \tag{F.17.2}
\]
The set $\Delta$ consists of sixteen distinct points of $Y$.
Set
\[
 Y^\dagger=Y\setminus\Delta,\qquad
 Y_0=\mathbf P^1\setminus\{0,\infty,\zeta,-\zeta\}.
                                                               \tag{F.17.3}
\]
There is an actual ten-label configuration map
\[
 \phi^\dagger(s)=
       (0,\infty,1,(v_u(s))_{u\in\mu_5},a(s),-a(s))
             \quad(s\in Y^\dagger).
                                                               \tag{F.17.4}
\]
The two added labels are ordered by these last two formulas.
At the basepoint they are the two fixed points of the actual
involution in F.12--15.  This is a specified ordering; comparison
with the other ordering is the literal transposition of these
two new labels.

The source fillings and their completed homomorphisms are
\[
 \widehat{\pi_1(Y^\dagger,s_0)}=\widehat F_{27}
       \xrightarrow{\ q\ }\widehat{\pi_1(Y,s_0)}=\widehat F_{11}
       \xrightarrow{\ q_0\ }\widehat{\pi_1(Y_0,s_0)}=\widehat F_3 .
                                                               \tag{F.17.5}
\]
The first fills precisely $\Delta$, and the second fills
the eight points $\pm u$ with $u\ne\zeta$.
Both maps are onto.

Each based meridian $m_d$ at an added point $d\in\Delta$
is killed by $q$, but
$\widehat{\phi^\dagger}_*(m_d)$ generates a full closed
copy of $\widehat\Z$ in the ten-point pure completion.
It is the actual primitive two-point twist at the corresponding
collision in (F.17.2), with the base path induced by the
specified source whisker.
Consequently the particular map
$\widehat{\phi^\dagger}_*$ does not factor through $q$.

Forgetting the four labels $v_u$ with $u\ne\zeta$ gives
the exact map
\[
 \chi:Y_0\longrightarrow\mathcal M_{0,6},\qquad
                 \chi(s)=(0,\infty,1,a(s)^2,a(s),-a(s)).
                                                               \tag{F.17.6}
\]
Its induced full profinite map is injective and its discrete
image has the full induced topology.
The retained involution is $V\mapsto a(s)^2/V$.
Its actual branched quotient has the four labelled values
$0,\infty,1,t$, where $t=s^2/\zeta^2$.
The two maps $Y_0\to\mathcal M_{0,4}$ involved here are
specified in the proof: ordinary forgetting has parameter
$a^2$, whereas the involution quotient has parameter $t$.
Each is a degree-two unramified covering.
The first proves the full topology assertion directly.

These statements give an exact additional labelled
correspondence for the source-membership problem.
They make no injectivity assertion for the whole
ten-point map (F.17.4), and do not determine the
intersection in F.12.5.
\end{claim}

\begin{proof}
\emph{All collision points of the two new labels.}
The rational labels $a,-a$ are the fixed points of
$V\mapsto v_\zeta(s)/V$.
They are distinct and avoid $0,\infty,1,v_\zeta$
on $Y$: the excluded cases are $s=0,\infty,\zeta,-\zeta$.
For $u\ne\zeta$, the two exact differences have the
factorizations
\[
\begin{split}
 v_u-a
   &=\frac{2s(s^2-A_u)}{(u+s)^2(\zeta+s)},\\
 v_u+a
   &=\frac{2(\zeta-2u)(s^2-B_u)}
                {(u+s)^2(\zeta+s)}.
\end{split}                                                     \tag{F.17.7}
\]
This proves the stated collision equations, without
discarding a denominator or a source puncture.

Here is also a direct separation check.
In the specified complex embedding $|u|=|\zeta|=1$,
\[
 |A_u|=|2\zeta-u|>1,\qquad
                  B_u=\frac1{\overline{A_u}},\qquad |B_u|<1.
                                                               \tag{F.17.8}
\]
Equality in the first inequality would force $u=\zeta$.
For two distinct fifth roots $u,w$,
$A_u=A_w$ would imply $u+w=2\zeta$; the equality case
in the triangle inequality again forces $u=w=\zeta$.
Thus the four $A_u$ are distinct.  Their reciprocal
conjugates are distinct, and the two sets cannot
intersect by (F.17.8).
All eight values are nonzero.
Their square roots are therefore sixteen distinct
points.  None has absolute value one or equals
$0,\infty$, so none is an old source puncture.

The denominators in (F.17.7) are nonzero at these
sixteen points, since the excluded zeros have
absolute value one.
Every zero displayed there is simple:
the $s^2-A_u$ and $s^2-B_u$ factors have nonzero
derivatives at their roots, and the remaining
numerator factors are nonzero.
There is only one new collision at each such point.
The original eight labels are distinct on $Y$ by F.12.
This proves that (F.17.4) is defined precisely on
$Y^\dagger$ among the points of $Y$.

The basepoint $s_0$ belongs to $Y^\dagger$.
Indeed its actual involution permutes all eight
original labels in four nontrivial pairs, by F.15.10.
Neither of its two fixed points is one of those labels.
The formula in (F.17.1) squares to $v_\zeta(s_0)=v_2(s_0)$,
so it selects exactly the required ordered new pair.

\emph{The actual filling maps and local twists.}
Choose a fixed system of based geometric whiskers
from $s_0$ to small positively oriented circles at
the finite punctures of $Y^\dagger$.
The circles are disjoint, and the whisker interiors
avoid all punctures.
The corresponding meridians are a free basis, with
the infinity meridian given by their inverse ordered
product.
Filling the sixteen punctures $\Delta$ sends precisely
those meridians to the identity and retains the
meridians of $Y$ with their induced whiskers.
Filling the eight additional old punctures gives $Y_0$.
The same maps after completion are exactly (F.17.5):
the finite-target universal property identifies
each filling kernel with the closed normal closure
of its named meridians.  In particular both completed
maps are onto, and their composite fills all twenty-four
additional points of $Y^\dagger$ relative to $Y_0$.

Let $d^2=A_u$.
In a small disk about $d$, the difference $v_u-a$
has winding one by (F.17.7); every other label
difference stays nonzero there.
The local monodromy in the ordered configuration
is therefore the one-turn pure twist around the
two labels $v_u,a$.
Concretely, the center and difference of these two
coordinates give local configuration coordinates;
the center varies in a contractible disk, and the
nonzero difference runs once around zero.
Extending this representative by the identity outside
a small disk containing the two labels gives that
twist.  Its based conjugating path is the image
under $\phi^\dagger$ of the already chosen whisker.
For $d^2=B_u$ the identical argument uses $v_u+a$
and the pair $v_u,-a$.
Thus the asserted monodromy includes the actual
pair of labels, exponent, and base transport.

For a full profinite detection, normalize the original
three labels $0,\infty,1$ in the ten-point moduli space.
Write $Z_u,Z_+,Z_-$ for the coordinates of $v_u,a,-a$.
The globally nonzero functions $Z_u-Z_+$ and $Z_u-Z_-$
define their integer winding homomorphisms on its
pure mapping class group.
Their pullbacks under $\phi^\dagger$ are precisely
the two functions in (F.17.7).
At the corresponding added meridian the relevant
integer winding is one.  Completing that character
gives the literal equality
\[
 \widehat{\operatorname{wind}}_{u,\pm}
       \bigl(\widehat{\phi^\dagger}_*(m_d^{\,z})\bigr)
                         =z\quad(z\in\widehat\Z).
                                                               \tag{F.17.9}
\]
The continuous power map onto the closure of that
twist is injective by (F.17.9), as well as surjective.
This proves the full $\widehat\Z$ assertion at every
primary coordinate.  It does not require injectivity
of $\widehat{\phi^\dagger}_*$ on the entire source.

Since $q(m_d)=1$ but its ten-point image is nonidentity,
there cannot be a homomorphism $D$ with
$Dq=\widehat{\phi^\dagger}_*$.
This proves exactly the stated nonfactorization.

\emph{Both diagrams with their literal retained labels.}
Let $\operatorname{forget}_{+,-}$ forget the two added
labels, and let $\operatorname{forget}_{\mathrm{four}}$
forget precisely the four $v_u$ with $u\ne\zeta$.
The based maps on full completions satisfy
\[
\begin{split}
 \operatorname{forget}_{+,-}\,
                    \widehat{\phi^\dagger}_*
       &=\widehat\phi_*\,q,\\
 \operatorname{forget}_{\mathrm{four}}\,
                    \widehat{\phi^\dagger}_*
       &=\widehat\chi_*\,q_0q .
\end{split}                                                     \tag{F.17.10}
\]
The first target has the original eight labels.
The second target has exactly
$0,\infty,1,v_\zeta,a,-a$.
Both identities follow from the actual configuration
coordinates, with the same basepoint and whiskers.
Their common further forgetting retains
$0,\infty,1,v_\zeta$ and has parameter
\[
                          u_0(s)=a(s)^2.                        \tag{F.17.11}
\]
This common four-point target is the ordinary
forgetting target, not the branched quotient target
described next.

\emph{The six-point double cover and its full topology.}
The coordinate $a=(\zeta-s)/(\zeta+s)$ is a M\"obius
isomorphism from $Y_0$ onto
$\mathbf P^1\setminus\{0,\infty,1,-1\}$.
Squaring $a$ is the degree-two unramified covering
of $\mathbf P^1\setminus\{0,\infty,1\}$.
It is exactly the ordinary forgetting (F.17.11).
The induced fundamental-group image is the index-two
kernel of the character that is one on a meridian
at zero and zero on a meridian at one, using the
$u_0$ coordinate and its induced basepoint.

A finite-index subgroup of a discrete group has
its full induced profinite topology: the normal core
of a finite-index subgroup of it has finite index
in the whole group.
Consequently the completed map on this covering is
an embedding $\widehat F_3\hookrightarrow\widehat F_2$.
Since it is the composite of $\widehat\chi_*$
with an actual target forgetting homomorphism,
$\widehat\chi_*$ is injective.
The same composite proves the full induced topology
of the discrete image of $\chi_*$.
This proof asserts full topology for this specified
six-point map and uses no inference about the kernel
of the ten-point map.

The involution $V\mapsto a^2/V$ fixes the two labels
$a,-a$ individually and exchanges the other labels
as $(0\ \infty)(1\ a^2)$.
Its actual quotient is
\[
 Q_a(V)=\left(\frac{V-a}{V+a}\right)^2,
 \qquad
 \begin{array}{c|c}
 \text{labelled fibre}&\text{quotient value}\\ \hline
 a&0\\
 -a&\infty\\
 \{0,\infty\}&1\\
 \{1,a^2\}&t=s^2/\zeta^2 .
 \end{array}                                                    \tag{F.17.12}
\]
Direct substitution proves all four entries;
in particular
$Q_a(1)=((1-a)/(1+a))^2=s^2/\zeta^2$.
Thus the branched-quotient parameter map is the
other degree-two unramified covering
$Y_0\to\mathbf P^1\setminus\{0,\infty,1\}$.

The sheet change $s\mapsto-s$ sends $a$ to $a^{-1}$.
The M\"obius map $V\mapsto V/a^2$ identifies the
two six-point covers over the same $t$ and exchanges
exactly the labels $1,a^2$, while preserving the
two named branch labels and the labels $0,\infty$.
The deck involution exchanges both ordinary pairs.
Hence that sheet change cannot be corrected to a
pure labelling by the deck involution.
Conversely the choices of the two ordinary
sheet orderings, modulo the simultaneous deck
exchange, are exactly these two possibilities.
This proves that $Y_0$ is the actual index-two
pure-labelling cover of this four-point quotient
family, with all of its labels and the branch order
specified.

The executed certificate
\path{code/part_a/common_reflection_drilling_f17.py}
checks the eight exact factorizations over
$\mathbf Q(\zeta)$, all separations of the eight
squared collision values, the two quotient and
sheet-label formulas, and the coefficient-one
meridian detectors in finite primary reductions.
It passes 445 exact checks.
The topology and full completion comparisons are
the explicit covering and winding arguments above.
\end{proof}
